\chapter{Latency, reliability and performance}
\label{ch:performance}

\fxme\ trades latency for reliability. This chapter explains the two
settings that control the trade, how to read what happens, and how to tune
the machines.

\section{Block size N and network budget D}

The plugin sends the audio to the server in blocks of N samples and gives
each reply D blocks of time to come back. Both are in the \ui{Settings}
tab.

\begin{description}
  \item[Block size N] 32 to 1024 samples. A small block means less latency
    but more packets per second and tighter deadlines for the server and
    the network.
  \item[Network budget D] 1 to 8 blocks. A larger budget survives a slower
    network or a busier server, at the price of more latency.
\end{description}

The latency that \fxme\ reports to the DAW is roughly N times D samples
plus the DAW's own block size, plus the latency of the chain's plugins. The
\ui{Settings} tab shows the exact figure for your settings, in samples and
milliseconds, split between the network and re-blocking part and the chain
itself. The DAW compensates for all of it, so the track stays in time; what
changes is how late a played instrument is heard, which matters when you
record through the plugin.

If a reply is not back in time, that block is \emph{concealed}: a short fade
to silence (64 samples) and back, so that there is no click, only a
very brief dip. Concealed blocks are counted on the bottom line of the
window.

The default is N\,=\,128 and D\,=\,2 (5.8\,ms of budget at 44.1\,kHz).

\section{Choosing the settings}

\begin{longtable}{@{}p{0.3\textwidth}p{0.64\textwidth}@{}}
\toprule
\textbf{Situation} & \textbf{Suggested settings}\\
\midrule
\endhead
Mixing, many heavy tracks & 1024$\times$2 or 512$\times$2: the most robust. Latency is no concern while mixing.\\
Recording through the plugin, light chain & 128$\times$2, or 64$\times$2 if the server has a lot of spare CPU.\\
A server near full load & Larger blocks. A busy server sometimes stalls for 20 to 50\,ms, which small blocks cannot absorb.\\
Network with other traffic or Wi-Fi & Raise D, or move to a wired connection.\\
\bottomrule
\end{longtable}

\begin{tip}
Test with the DAW's own settings in mind: the DAW's audio buffer adds to
\fxme's latency. Using a large buffer in the DAW for mixing and a small one
for recording is common practice, and works the same here.
\end{tip}

\section{Measured performance}

These figures come from the first stress tests (details in the project's
\file{doc/Tests} folder), on old machines: a 2015 Intel Core i7-6700K for the
client (Windows, Reaper, an RME ASIO interface at 44.1\,kHz and a buffer of
128) and a 2018 AMD Ryzen 5 PRO 2400G for the server (Ubuntu 24.04, four
cores). Every track had one \ui{FxmeNet} instance, and the ``full chain''
holds all 17 plugins of the FxmeFX collection (FX-Mechanics' effects).

\begin{longtable}{@{}p{0.22\textwidth}p{0.12\textwidth}p{0.14\textwidth}p{0.4\textwidth}@{}}
\toprule
\textbf{Tracks} & \textbf{Chain} & \textbf{N$\times$D} & \textbf{Result}\\
\midrule
\endhead
24 local & full & (local) & 75\,\% of the client's CPU.\\
24 remote & full & 256$\times$2 & Client at 25\,\%, server at about 66\,\%, about 10 concealed blocks in 5 minutes on one track.\\
24 remote & full & 1024$\times$2 & Client under 20\,\%, almost no concealed blocks.\\
28 remote & full & 1024$\times$2 & Server above 75\,\%: no concealed blocks, or very rare ones. Very usable.\\
32 remote & one equaliser & 128$\times$2 & Network at 95\,Mbit/s each way, no problem.\\
\bottomrule
\end{longtable}

What these tests show:

\begin{itemize}
  \item The offload works: the client's processor load fell from 75\,\% to
    20 to 25\,\% for the same session.
  \item Heavy loads (28 tracks with 17 plugins each, 476 plugin instances)
    are very usable with large blocks.
  \item Small blocks need headroom. While the server has spare CPU, 128 and
    256 work well; close to full load the server stalls for a few tens of
    milliseconds now and then, which short deadlines cannot absorb.
  \item These are extreme loads. Real sessions rarely put 17 plugins on every
    track.
\end{itemize}

\section{Making the server reliable (Linux)}
\label{sec:linuxtuning}

A generic Linux system works. These are the usual steps for audio on Linux;
try them one at a time and check after each whether the concealed count goes
down. More testing will be done. Feedback is welcome on the performance you get on other systems and with other optimizations.

\begin{enumerate}
  \item \textbf{Real-time priority for the audio threads.} Check with
\begin{lstlisting}
ps -eLo tid,cls,rtprio,pcpu,comm | grep -i fxme
\end{lstlisting}
    while a session is open: the audio threads should be in class \cmd{RR}
    or \cmd{FF}. If they are \cmd{TS}, the system refused real-time priority
    (the server's log says so). Create
    \file{/etc/security/limits.d/audio.conf}:
\begin{lstlisting}[language={}]
@audio - rtprio 95
@audio - memlock unlimited
\end{lstlisting}
    Add the account that runs the server to the \cmd{audio} group, and log in
    again (or restart the service).
  \item \textbf{Bigger socket buffers.} Under heavy load the kernel can drop
    replies when the send buffer is full (the server's log counts them as
    \emph{send errors}; \cmd{nstat -az UdpSndbufErrors} shows the
    system-wide count). Raise the limits, then restart the server:
\begin{lstlisting}
sudo sysctl -w net.core.rmem_default=4194304 net.core.rmem_max=8388608
sudo sysctl -w net.core.wmem_default=4194304 net.core.wmem_max=8388608
\end{lstlisting}
    To keep them after a reboot, put the four lines (without \cmd{sudo
    sysctl -w}) in \file{/etc/sysctl.d/90-fxmenet.conf}.
  \item \textbf{CPU frequency.} Set the governor to performance:
\begin{lstlisting}
sudo cpupower frequency-set -g performance
\end{lstlisting}
  \item \textbf{Kernel.} A low-latency or fully preemptible kernel
    (\cmd{linux-lowlatency}, or booting with \cmd{preempt=full}) helps the
    worst cases.
\end{enumerate}

\todo{When the server asks for large socket buffers itself (todo item 21),
step 2 becomes: ``the log tells you if the system limits are too low''.}

\section{Making the client reliable}

\begin{itemize}
  \item Use a wired gigabit connection, and keep other heavy traffic off it.
  \item In Reaper, try turning off \emph{anticipative FX processing} for the
    tracks that hold \ui{FxmeNet}, and compare the concealed counts.
  \item On the network adapter, turn interrupt moderation off or to low, and
    use the high-performance power plan.
\end{itemize}
